\documentclass[11pt]{article}
\usepackage[margin=0.85in]{geometry}
\usepackage{amsmath,amssymb,amsthm,booktabs,longtable,array,hyperref}
\setlength{\emergencystretch}{2em}
\hypersetup{colorlinks=true,linkcolor=blue,urlcolor=blue,citecolor=blue}
\newtheorem{theorem}{Theorem}[section]
\newtheorem{lemma}[theorem]{Lemma}
\newtheorem{proposition}[theorem]{Proposition}
\newcommand{\F}{\mathbb F_2}
\newcommand{\Z}{\mathbb Z}
\newcommand{\Nf}{\mathsf N_{4,0}}
\newcommand{\repo}{https://github.com/alectionary/vibe-math}
\newcommand{\prooflink}[2]{\href{https://github.com/alectionary/vibe-math/blob/41307a70c2c0952e2222ee20ae4517a9b046030f/formalization/lift-label-elimination/#1.lean}{\texttt{#2}}}
\newcommand{\priorproof}[2]{\href{https://github.com/alectionary/vibe-math/blob/41307a70c2c0952e2222ee20ae4517a9b046030f/formalization/absorbing-four-zero/#1.lean}{#2}}
\title{Cyclic-voltage elimination of Pasch supports\\in two cubic configurations}
\author{Project Contributors}
\date{24 October 2026}
\begin{document}
\maketitle
\begin{abstract}
We prove exact elimination criteria and complete label counts for elementary
$(4,0)$ absorbing supports in cyclic covers of two cubic, linear $10_3$
configurations, and construct a binary CSS family with no such supports in
either Tanner graph. The problem follows the four-variable counting direction
identified by Abiad and Castriota (15 July 2026, Remark 4.18 and Section 6).
For an ordinary $N$-fold monomial cover,
every elementary $(4,0)$ support projects injectively to a base Pasch and
contributes exactly $N$ supports precisely when its three fundamental cycle
voltages vanish. Explicit unit normal forms over the integers give complete
label histograms for every $N\geq1$, including composite $N$. The numbers of
eliminating raw label arrays are $N^{30}-N^{27}$ and
$N^{30}-5N^{27}+10N^{24}-10N^{21}+4N^{19}$.
We distinguish elimination in a cover sector from elimination in the full
scalar-right CSS construction. The fixed seven-coordinate scalar forces at
least seventy supports per Tanner graph. Replacing its length by $7q$ and
changing just two base-edge labels yields a family with parameters
$[[140q,18,6]]$, $q\geq2$, column weight three, row weight six and no
$(4,0)$ supports in either graph. Every proof has a checked Lean version.
\end{abstract}
\noindent\textbf{Verification.} Every manuscript proof has a checked Lean version. The 70 mathematical modules were built from their recorded sources; the axiom audit covers all 1,755 named authored declarations. The source links below are pinned to \href{https://github.com/alectionary/vibe-math/commit/41307a70c2c0952e2222ee20ae4517a9b046030f}{\texttt{41307a70c2c0}}. The machine-readable \href{https://github.com/alectionary/vibe-math/blob/41307a70c2c0952e2222ee20ae4517a9b046030f/formalization/lift-label-elimination/proof-coverage.json}{proof inventory} records the claim-to-theorem correspondence; \href{https://github.com/alectionary/vibe-math/blob/41307a70c2c0952e2222ee20ae4517a9b046030f/formalization/lift-label-elimination/PROOFS.md}{the full file index} links every compiled proof and preserved research import.

\section{Questions and answers}
Abiad and Castriota~\cite{AC}, in the first version of their preprint dated
15 July 2026, identify counting elementary $(4,2)$ and $(4,0)$ absorbing sets
as further work. Remark 4.18 defers the $(4,2)$ computation and leaves the
$(4,0)$ case undeveloped; the final bullet of Section 6 proposes accessing
both through the same spectral data. Their question led to our earlier
incidence-corrected counting formulas and counterexamples to sufficiency of
the evaluated matrices $M_j$. This paper takes the next step: it determines
exactly which labels eliminate the $(4,0)$ supports in two fixed bases,
counts every such label assignment, and proves an explicit CSS construction.

The following table separates the source-posed counting direction, answered
in our preceding work, from the elimination questions answered here. The
source date precedes this revision by 101 days.
\begin{longtable}{@{}>{\raggedright\arraybackslash}p{0.28\textwidth}>{\raggedright\arraybackslash}p{0.27\textwidth}>{\raggedright\arraybackslash}p{0.39\textwidth}@{}}
\toprule Question&Source and date&Specific answer and proof\\\midrule\endhead
How can the $(4,2)$ and $(4,0)$ supports be counted? Can the $M_j$ data alone determine them?&
Abiad--Castriota, 15 July 2026, Remark 4.18 and Section 6.&
Earlier work: exact incidence-corrected formulas for both counts; a negative
answer to $M_j$ sufficiency. In particular, two $(3,6)$-regular 140-variable
examples have identical full $M_j$ at every frequency but $(4,0)$ counts
77 and 105. \priorproof{AbsorbingFourCounts}{Combined formulas};
\priorproof{PaschCubicTrade}{counterexample}.\\[4pt]
Which cyclic labels eliminate every $(4,0)$ support in the two bases, and how many are there?&
Our elimination follow-up to that counting direction, 24 October 2026.&
Three necessary and sufficient cycle equations per base support
(Theorem~\ref{thm:voltage}); exact counts for every $N\geq1$
(Theorem~\ref{thm:hist}); simultaneous original/dual count
(Proposition~\ref{prop:joint}).\\[4pt]
Does eliminating cover supports eliminate them in the full CSS graphs?&
Our full-incidence follow-up, 24 October 2026.&
Theorem~\ref{thm:full} gives an exact sector partition. At scalar length seven,
at least 70 supports remain per full graph, whatever the cover labels.\\[4pt]
Can both full graphs be free of these supports while retaining a nonzero quantum code?&
Our construction question, 24 October 2026.&
Yes. Theorem~\ref{thm:family} constructs $[[140q,18,6]]$ codes for every
$q\geq2$, with column weight three, row weight six and no elementary
$(4,0)$ support in either graph.\\
\bottomrule
\end{longtable}
The source's spectral proposal is a research direction, not a numbered
conjecture. The counterexample disproves its $M_j$-only interpretation; the
corrected formulas answer the counting question using the additional incidence
information. The elimination questions in the last three rows are ours.

\subsection*{Main results}
For every positive lift size $N$, we prove that the numbers of raw edge-label
arrays eliminating all cover supports are
\[
 E_{\rm star}(N)=N^{30}-N^{27},\qquad
 E_{\rm triangle}(N)=N^{30}-5N^{27}+10N^{24}-10N^{21}+4N^{19}.
\]
The complete histograms hold for composite and even $N$ as well as prime $N$.
We then prove why eliminating the cover sector alone fails at scalar length
seven, and overcome that obstruction by using length $7q$ and changing exactly
two base-edge labels. The resulting family has exact dimension 18 and exact
mixed-Pauli distance six for every $q\geq2$. Sections 2--4 prove these results;
Section 5 links each proof to its checked Lean sources.

\subsection*{Conventions}

A variable neighborhood is a finite set $H(v)\subseteq C$. We assume
$|H(v)|=3$ and $|H(u)\cap H(v)|\leq1$ for $u\ne v$. An elementary $(4,0)$
support is an \emph{unordered} set of four variables for which every used
check has selected degree two. In a cubic linear incidence structure it is a
Pasch: the normal graph is $K_4$, with six distinct pair checks. Equivalently,
its four actual binary columns have zero sum. A selected degree-four check
would consume the unique common check of every pair and prevent pairing the
remaining sockets, so a four-column binary dependency is automatically
elementary under these hypotheses. We use the preserved generic Lean
classification and actual binary-matrix bridge for these facts.

For CSS parameters we use the binary check-matrix model: $XZ^{\mathsf T}=0$,
$k=n-\operatorname{rank}X-\operatorname{rank}Z$, and distance equal to the
minimum union support of a nontrivial logical pair $(x,z)$. Such a pair obeys
$Zx=Xz=0$ and lies outside the product of the two stabilizer row spaces.
The Lean statements use these actual matrices and this mixed-Pauli convention.

\subsection*{Related work}
LDPC constructions based on voltage graphs were studied by Kelley and
Walker~\cite{KW}. Asvadi et al.~\cite{ABA} explicitly relate cyclic
voltage sums to lifted cycle closure and intentional removal of trapping
structures. The cycle-consistency matrix approach of
Wang et al.~\cite{WDW11,WDW12} gives the corresponding fundamental-cycle
framework for absorbing structures. Its SCB separability, prime-field and
consistency hypotheses must be compared with the present arbitrary monomial
labels over $\Z/N\Z$; we prove that translation rather than treating the
settings as identical. Earlier combinatorial QC constructions and
Pasch-based distance arguments also appear in Vasi\'c and Djordjevi\'c~\cite{VD} and
Noda~\cite{Noda}.

The two base configurations below are classical; their construction and author
data are credited to~\cite{PP,AFKL}. Relevant passages and the author data were
compared in the preceding research block, with exact relabelings retained in
its artifact. The contributions here are the all-$N$ label counts, the
full-incidence obstruction and the explicit family with proved parameters.

The prior-art review covered classical literature, public GitHub repositories
and code, and authenticated X searches. The prior-art ledger records queries,
inspected matches and access details; Appendix A gives source dates and
inspected versions. The general voltage method and the base configurations
are established ingredients; the present theorems give the exact counts and
construction stated above.

\section{Projection and the exact voltage criterion}
Let $G$ be a finite abelian group and put a voltage $z(v,c)\in G$ on every
existing incidence. Define the actual cover columns by
\begin{equation}\label{eq:cover}
 H_z(v,t)=\{(c,t+z(v,c)):c\in H(v)\}.
\end{equation}
Voltages on nonedges have no effect. The numerical counts below specialize to
$G=\Z/N\Z$, $N\geq1$.

\begin{lemma}[Faithful Pasch projection]\label{lem:projection}
Every Pasch in $H_z$ has four distinct base variable types and projects to a
base Pasch.
\end{lemma}
\begin{proof}
Two different fibers of the same variable cannot share a check: their single
edges to a fixed base check would require $t+z=s+z$, hence $t=s$. The four
pairwise adjacent variables of a lifted Pasch therefore have distinct types.
Suppose three projected types contained the same base check $c$. Base linearity
forces their pair checks to have type $c$. The single $c$-edge at one lifted
variable then forces a common actual check of selected degree at least three,
a contradiction. Thus projected selected degrees are at most two. Cubicity and
pair completeness now give the base Pasch.
\end{proof}

Fix an ordered representative $v_0,v_1,v_2,v_3$ of a base Pasch, and let
$c_{ij}$ be its unique pair check. Set $d_{ij}=z(v_i,c_{ij})-z(v_j,c_{ij})$.
The six actual lifted pair incidences are the equations $x_j-x_i=d_{ij}$.
\begin{theorem}[Voltage criterion and unordered count]\label{thm:voltage}
The representative lifts if and only if
\begin{equation}\label{eq:balance}
 d_{01}+d_{12}=d_{02},\qquad d_{01}+d_{13}=d_{03},\qquad
 d_{02}+d_{23}=d_{03}.
\end{equation}
A balanced base support has exactly $|G|$ lifted supports. Consequently
\begin{equation}\label{eq:global}
 \Nf(H_z)=|G|\,\#\{\text{balanced base Pasch supports}\}.
\end{equation}
\end{theorem}
\begin{proof}
Cancellation gives necessity. Under~\eqref{eq:balance}, all six equations hold
for, and only for,
$(x_0,x_1,x_2,x_3)=(t,t+d_{01},t+d_{02},t+d_{03})$, $t\in G$.
The base pair checks are distinct, so these equations give an actual Pasch.
Distinct $t$ give distinct unordered supports because their base types differ;
different base supports have different projections. Lemma~\ref{lem:projection}
gives surjectivity. Formally, the last step is a bijection between unordered
supports with injective projection and a base support together with offsets
above a fixed labeling. Reindexing that labeling changes neither existence nor
multiplicity. The checked proof includes this finite-set partition, not just
the solution count for a fixed ordered tuple.
\end{proof}

\section{Complete label counts for the two fixed bases}
Write a triple such as $012$ for $\{0,1,2\}$. The star base has columns
\[
012,034,135,245,067,168,269,378,479,589.
\]
The triangle base replaces its first four triples by $345,125,024,013$.
Both have ten variables, ten checks and thirty edges, and are cubic and linear
on both sides. The star has the sole Pasch $0123$. The triangle has exactly
\[
0123,\quad0789,\quad1569,\quad2468,\quad3457.
\]
These complete finite lists and all pair checks are kernel-checked. A raw label
array has thirty coordinates, ordered by variable and then check. For the auxiliary eleven-coordinate slice, the allowed edge positions are
$7,10,11,16,19,22,23,25,26,28,29$ for the star and
$8,10,11,16,19,22,23,25,26,28,29$ for the triangle (zero-based indices);
the remaining nineteen labels are zero. The label encodings are injective.
We count raw arrays and these explicit slices independently. No assertion
about equivalence classes under switching is needed or made.

For a subset $J$ of the balance events, stack its three integer cycle rows
into $A_J$. Explicit elementary integer row and column operations reduce
$A_J$, both on all thirty edges and on the eleven slice coordinates, to
$\operatorname{diag}(I_r,0)$. For the star $r=0,3$. For the triangle,
\begin{center}
\begin{tabular}{c@{\qquad}rrrrrr}
\toprule
$|J|$&0&1&2&3&4&5\\
$r_J$&0&3&6&9&11&11\\
\bottomrule
\end{tabular}
\end{center}
The nonzero invariant factors are all units. This integer assertion is stronger
than a rational or prime-field rank computation.

\begin{lemma}[Certificate soundness]\label{lem:integer}
If integer matrices $U,U^{-1},V,V^{-1}$ satisfy the recorded inverse identities
and $UAV=\operatorname{diag}(I_r,0)$, reduction modulo every positive $N$
gives $N^{n-r}$ solutions to $Az=0$ in $(\Z/N\Z)^n$.
\end{lemma}
\begin{proof}
The integer identities survive the ring homomorphism to $\Z/N\Z$. The map
$z\mapsto V^{-1}z$ is a bijection from the original kernel to the diagonal
kernel, with inverse $y\mapsto Vy$. The diagonal kernel consists of vectors
with its first $r$ coordinates zero and its remaining $n-r$ arbitrary. This
argument needs neither a field nor coprimality. All inverse and normal-form
identities, as well as their connection to the literal cycle rows, are checked
in Lean. An independent Python verifier reconstructs and replays the same
certificates.
\end{proof}

\begin{theorem}[All-$N$ label histograms]\label{thm:hist}
The raw numbers eliminating every left-cover Pasch are
\begin{align}
 E_{\rm star}(N)&=N^{30}-N^{27},\label{eq:star}\\
 E_{\rm triangle}(N)&=N^{30}-5N^{27}+10N^{24}-10N^{21}+4N^{19}.\label{eq:triangle}
\end{align}
For the triangle, let $g_j$ count slice arrays with exactly $j$ balanced
base Pasches. Then
\begin{align*}
 g_0&=N^{11}-5N^8+10N^5-10N^2+4,\\
 g_1&=5(N^8-4N^5+6N^2-3),\\
 g_2&=10(N^5-3N^2+2),\\
 g_3&=10(N^2-1),\qquad g_4=0,\qquad g_5=1.
\end{align*}
Raw histograms are $N^{19}g_j$. The star slice histogram is
$(N^{11}-N^8,N^8)$. An array in bin $j$ gives $Nj$ actual lifted supports.
\end{theorem}
\begin{proof}
Lemma~\ref{lem:integer} gives $N^{n-r_J}$ for every event intersection.
Inclusion-exclusion gives~\eqref{eq:star}--\eqref{eq:triangle}. More generally,
\[
 g_j=\sum_{s=j}^{5}(-1)^{s-j}\binom{s}{j}\binom{5}{s}N^{11-r_s}.
\]
This follows pointwise: a truth pattern with $k$ events contributes
$\sum_{s=j}^k(-1)^{s-j}\binom{s}{j}\binom{k}{s}$, which is one for $k=j$
and zero otherwise. Summing gives the histogram. The Lean proof checks the
finite truth-pattern identity, proves the summation bridge and groups subsets
by cardinality. It then evaluates the six polynomials. Literal voltage data
connect the abstract linear forms to the actual edge labels of both bases.
The raw calculation uses $n=30$, the slice calculation $n=11$; their ratio
and the displayed raw factors follow from the two independently checked counts.
\end{proof}

At $N=1$ no base Pasch can be eliminated. For $N\geq2$ the eliminating counts
are positive. In particular,
\[
 g_0=N^8(N^3-5)+10N^2(N^3-1)+4>0.
\]
For the star, $N^{30}-N^{27}=N^{27}(N^3-1)>0$.

\begin{proposition}[Simultaneous original and dual elimination]\label{prop:joint}
For the star base, the number of raw arrays eliminating both its original and
its dual-cover Pasch is
\begin{equation}\label{eq:joint}
 N^{30}-2N^{27}+N^{24}.
\end{equation}
\end{proposition}
\begin{proof}
The dual has its sole Pasch on check types $6,7,8,9$. Pulling its voltages back
to the original edges with a minus sign gives a second three-row event. The
stacked six-row matrix has an explicit unit normal form of rank six. Each
individual event has rank three. Lemma~\ref{lem:integer} and two-event
inclusion-exclusion give~\eqref{eq:joint}.
\end{proof}

\section{Full scalar-right incidence and the quantum family}
Let $R_N=\F[x]/(x^N+1)$, and let $A_{cv}=x^{z(v,c)}$ on each base edge.
The adjoint $*$ transposes a matrix and replaces $x$ by $x^{-1}$. For
$b=1+x+x^3$, use
\begin{equation}\label{eq:css}
 X=[A\mid b^*I_C],\qquad Z=[bI_V\mid A^*].
\end{equation}
Thus $XZ^*=Ab^*+b^*A=0$. Equivalently, common cyclic translations commute
with the arbitrary monomial cover, and the two equal binary summands cancel.
The Lean proof uses the actual expanded sparse matrices and block product.

An $X$ left column $(v,t)$ is~\eqref{eq:cover}; an $X$ right column $(c,t)$
has checks $(c,t-e)$ for $e\in\{0,1,3\}$. A $Z$ left column $(v,t)$ has
checks $(v,t+e)$, and a $Z$ right column $(c,t)$ has checks $(v,t-z(v,c))$.
These signs specify the ordinary binary transpose convention.

\begin{theorem}[Full support partition]\label{thm:full}
For $N\geq7$, any base whose incidence and transpose are cubic and linear, and any
monomial labels, suppose the scalar columns $B$ are cubic and linear. Then
\begin{equation}\label{eq:full}
 \Nf(X)=\Nf(H_z)+|C|\Nf(B),\qquad
 \Nf(Z)=\Nf(H^*_{z})+|V|\Nf(B).
\end{equation}
No mixed-sector Pasch occurs. For the two ten-type bases and $N=7$, each full
Tanner graph therefore retains at least seventy supports, for every left label
assignment.
\end{theorem}
\begin{proof}
An adjacent mixed pair has no proper common neighbor: base linearity and the
single edge of each left variable at a fixed check type would force any such
neighbor into a triple common check. In a Pasch every third variable would be
a proper common neighbor of either root pair, so a mixed pair is impossible.
Every support is all left or all right. Pairwise adjacency forces an all-right
support into a single check-type fiber. The finite-set partition gives the
sum of the left-cover count and one scalar count per fiber. The scalar at
$N=7$ has exactly seven supports; its transpose has the same count. Hence
$10\cdot7=70$ is unavoidable even when every left-cover Pasch is eliminated.
\end{proof}

\begin{theorem}[Explicit family]\label{thm:family}
For every integer $q\geq2$, put $N=7q$, use the star base and set
\begin{equation}\label{eq:labels}
 z(0,0)=z(4,6)=7,\qquad z(v,c)=0\text{ on all other edges}.
\end{equation}
The actual binary checks~\eqref{eq:css} define a CSS code with parameters
$[[140q,18,6]]$. Both graphs have column weight three, row weight six, no
Tanner four-cycle and no elementary $(4,0)$ support.
\end{theorem}
\begin{proof}
\emph{Scalar factor.} The kernel equation of the negative-shift scalar is
$f_{t+3}=f_t+f_{t+1}$. Every solution is determined by three bits $(a,b,c)$
and follows the seven-periodic pattern
\[
 (a,b,c,a+b,b+c,a+b+c,a+c).
\]
At length $7q$ all eight patterns occur. Every nonzero pattern has four ones
per seven coordinates, so every nonzero scalar kernel word has weight $4q$.
Thus $q\geq2$ excludes weight four. Reflection proves the same assertion for
the transposed scalar. The differences of $\{0,1,3\}$ are distinct modulo
$N\geq14$, proving scalar cubicity and linearity.

\emph{Cover sectors.} The original sole Pasch has a cycle residual $7$,
and the dual sole Pasch has a residual $\pm7$; these are nonzero modulo $7q$.
The checked duality maps base variable $u$ to dual variable $u+6\pmod{10}$
and base check $c$ to dual check $c+4\pmod{10}$, changing the two-edge label
pattern to its negative. Hence both cover sectors have no four-column binary
dependency. Theorem~\ref{thm:full} excludes supports in the full graphs.
Cubicity follows from the two three-socket column types. Cover linearity,
scalar linearity and the one-edge-per-check-type property exclude Tanner
four-cycles, including mixed pairs. A row has three cover neighbors and three
scalar neighbors, giving weight six.

\emph{Dimension.} In either check-equation kernel, the scalar conditions force
each coordinate sequence into a seven-periodic three-bit pattern (with
reflection for the negative orientation). All nonzero labels are multiples of
seven, so shifts by those labels act trivially on such sequences. Each of the
three phases therefore belongs to the kernel of the binary star incidence or
its transpose. Each base kernel has exactly eight elements, as certified by
its complete finite support list. Consequently each expanded check-equation
kernel has $8^3=512=2^9$ elements. Binary rank-nullity gives
\[
 \operatorname{rank}X=\operatorname{rank}Z=10N-9,
 \qquad k=20N-2(10N-9)=18.
\]
This recurrence argument also covers even $q$, without a square-free or
coprimality assumption on the polynomial quotient.

\emph{Distance.} Put ones on left variable types $\{2,3,5,6,7,8\}$ at fiber zero
and zero elsewhere, obtaining $w$ of weight six. These six base columns sum
to zero and avoid the two modified variables, so $Xw=0$. Let $p$ be the scalar
pattern with initial bits $(1,0,0)$, and put $u(v,t)=p(-t)$ for left type $v=2$,
with zero on all other coordinates. Then $Zu=0$ and $\langle w,u\rangle=1$.
Every vector in the row space of $Z$ is orthogonal to $u$, so $w$ is a nontrivial
logical word. Column weight three forces every kernel word to have even weight;
linearity excludes weight two; the elimination proof excludes weight four.
Thus both classical check kernels have minimum nonzero weight at least six.
Every nontrivial mixed CSS logical has a nontrivial component whose support
is contained in its Pauli support, so its weight is at least six. The pure word
$w$ attains six, proving the exact distance.
\end{proof}

The family gives arbitrarily long codes with both Tanner graphs free of the
specified supports and with exact dimension and distance. The case $q=2$ is
$[[280,18,6]]$. An independent C++ verifier reconstructs all
columns, checks both ranks $131$, commutation, degrees and linearity, and
exhaustively tests all $250{,}654{,}530$ four-subsets in each graph. It also checks
the logical support with zero-based indices $28,42,70,84,98,112$ by an augmented
row-rank calculation. Python exact consumers additionally cover $q=3,4$.
These computations are independent checks of the specified finite instances;
the all-$q$ conclusion uses the proof above.

\section{Proof artifacts and reproducibility}
The proof inventory below names checked entry points. All seven manuscript
proof blocks, plus their finite data, source translations and auxiliary claims,
have Lean versions with matching hypotheses and conclusions. The axiom audit
queries every named authored declaration, including definitions containing proofs. Source files and all imported research modules are
included in the repository. Compiled objects, third-party PDFs and local caches
are not redistributed.
\begin{longtable}{@{}>{\raggedright\arraybackslash}p{0.34\textwidth}>{\raggedright\arraybackslash}p{0.60\textwidth}@{}}
\toprule Claim or bridge&Direct repository proof file\\\midrule\endhead
Projection and actual incidence&\prooflink{PaschLiftProjection}{PaschLiftProjection.lean}; \prooflink{PaschLiftCriterion}{PaschLiftCriterion.lean}\\
Unordered all-$N$ count&\prooflink{SupportVoltageCount}{SupportVoltageCount.lean}; \prooflink{CountRepresentatives}{CountRepresentatives.lean}\\
Complete base lists and physical voltages&\prooflink{StarVoltageData}{StarVoltageData.lean}; \prooflink{TriangleVoltageData}{TriangleVoltageData.lean}\\
Integer certificate soundness&\prooflink{NormalFormCount}{NormalFormCount.lean}; individual normal forms and row connections in the manifest\\
Histograms and raw label counts&\prooflink{LabelHistograms}{LabelHistograms.lean}; \prooflink{ConcreteElimination}{ConcreteElimination.lean}\\
Label encodings and positivity&\prooflink{LabelEncoding}{LabelEncoding.lean}; \prooflink{PositiveLabelCounts}{PositiveLabelCounts.lean}\\
Simultaneous star/dual count&\prooflink{JointElimination}{JointElimination.lean}\\
Full scalar-right count&\prooflink{QCSupportCount}{QCSupportCount.lean}; \prooflink{FullGraphCounts}{FullGraphCounts.lean}\\
Fixed seven-coordinate obstruction&\prooflink{FanoObstruction}{FanoObstruction.lean}\\
Scalar factor, all $q$&\prooflink{ScalarWeight}{ScalarWeight.lean}; \prooflink{ScalarIncidence}{ScalarIncidence.lean}\\
Actual CSS matrices and dimensions&\prooflink{SparseCSSMatrices}{SparseCSSMatrices.lean}; \prooflink{FamilyRank}{FamilyRank.lean}\\
Full-graph elimination and distance&\prooflink{FamilyGraphs}{FamilyGraphs.lean}; \prooflink{FamilyPauliDistance}{FamilyPauliDistance.lean}\\
Row degrees and final family&\prooflink{SparseRowDegrees}{SparseRowDegrees.lean}; \prooflink{FamilyTheorem}{FamilyTheorem.lean}\\
\bottomrule
\end{longtable}
\noindent The pinned environment is Lean \texttt{v4.36.0-rc1}, mathlib revision
\href{https://github.com/leanprover-community/mathlib4/commit/5f7b4198d22a9c82e1c698a379334fe128095c82}{\texttt{5f7b4198d22a}} (full revision in the manifest). The final modules were
built from their recorded sources and audited for axioms; only \texttt{propext},
\texttt{Classical.choice} and \texttt{Quot.sound} are permitted. No admitted
proof, native-decision axiom or custom mathematical axiom is accepted.
The new checks distinguish newly compiled modules from the unchanged,
previously freshly checked sixty-one-module dependency artifact.

For example, the following are ASCII renderings of checked entry-point
statements. The first has $N>0$; the last two have $q>0$, and $F=\Z/2\Z$.
Here \texttt{Logical} uses the two actual binary syndromes and excludes the
corresponding stabilizer row spaces; \texttt{pauliWeight} counts the union of
the two component supports.
{\small\begin{verbatim}
theorem star_raw_count :
  (Nat.card {z : Fin 30 -> ZMod N //
    Eliminates StarVoltageData.H (StarVoltageData.labelsRaw z)}
    : Int) = (N : Int)^30 - (N : Int)^27

theorem dimension :
  Fintype.card (Qubits q) - (X q).rank - (Z q).rank = 18 := by
  have hx := rankX q
  have hz := rankZ q
  have hq := NeZero.ne q
  rw [length]
  omega

theorem distance_six (hq : 2 <= q) :
  And (forall x z : Qubits q -> F,
         Logical q x z -> 6 <= pauliWeight q x z)
      (And (Logical q 0 (word q))
           (pauliWeight q 0 (word q) = 6))
\end{verbatim}}
The source proofs respectively transport the exact physical elimination
predicate to its counted balance event, apply binary rank-nullity to the
literal check matrices, and combine a lower bound for every mixed Pauli word
with the explicit separating logical witness. The complete statements,
including their section parameters and imports, are in the linked files.


Run \texttt{verify.py} for the independent certificate and C++ checks, and
\texttt{verify\_lean.py --project /path/to/pinned/project} after rebuilding the
preserved dependency artifact. The independent slice enumeration covers all
eleven-coordinate arrays for $N=2,3,4$ in both bases: $8{,}746{,}998$ arrays.
The polynomial histogram agrees in every case. The 128 seeded direct lifted
incidence consumers are sampled checks, not exhaustive evidence. Frozen
predecessor archives remain unchanged.

\section{Conclusions}
Theorem~\ref{thm:hist} completely answers the fixed-base label-counting
question for every positive lift size. Theorem~\ref{thm:full} identifies the
scalar obstruction that survives elimination in the cover sector.
Theorem~\ref{thm:family} overcomes it with two nonzero edge labels and proves
the exact parameters $[[140q,18,6]]$ for every $q\geq2$. Together these results
turn the four-variable counting investigation into an explicit construction
with both Tanner graphs free of elementary $(4,0)$ supports. Every manuscript
proof is linked to its checked Lean version. The constant dimension and
distance describe the family's parameters; decoding performance is a separate
experimental question.

\appendix
\section{Reference chronology and timing}
Every delta is the age of the stated source milestone at this paper's revision,
24 October 2026. It is \emph{not} a claim that a problem remained unsolved for
that interval. Year-only dates give a day-range; conference intervals are
retained. For the 2016 source we use the earliest located publisher date, not
an asserted earliest preprint. Date arithmetic is reproduced in
\texttt{reference-chronology.json}.
\begin{longtable}{@{}>{\raggedright\arraybackslash}p{0.19\textwidth}>{\raggedright\arraybackslash}p{0.20\textwidth}>{\raggedright\arraybackslash}p{0.39\textwidth}r@{}}
\toprule Source&Earliest located milestone&Inspected version / access&Delta (days)\\\midrule\endhead
Vasi\'c--Djordjevi\'c&1--3 Oct. 2003&Conference interval; primary institutional abstract only&8422--8424\\
Noda&3--7 Apr. 2006&Conference interval; primary publisher abstract only&7505--7509\\
Kelley--Walker&2008 (year only)&ISIT year; institutional abstract only&6506--6871\\
Asvadi et al.&23 Feb. 2010&arXiv v1 read; final 2011 metadata only&6087\\
Wang et al.&2011 (year only)&ITA author PDF read&5411--5775\\
Wang et al.&30 Aug. 2012&5 Dec. 2012 v2 read; final 2013 metadata only&5168\\
Pra\.{z}mowska and Pra\.{z}mowski&16 Feb. 2016&Publisher date; relevant HTML proofs read&3903\\
Abreu et al.&10 Apr. 2021&12 Sep. 2021 v2 passages and author data&2023\\
Abiad--Castriota&15 July 2026&v1 source conventions rechecked&101\\

\bottomrule
\end{longtable}
The bounded mathematical research block ran from 15:54:07 to 16:54:07 UTC on
24 October 2026 (sixty minutes). It included the prior-art audit and ordinary
construction/count arguments. Subsequent Lean proof transfer, manuscript
preparation and publication are recorded separately; that later work is not
presented as part of the sixty-minute investigation. Final publication timing
and commit identifiers are recorded in the artifact's publication manifest.

\begin{thebibliography}{9}\small
\bibitem{VD} B. Vasi\'c and I. B. Djordjevi\'c, Quasicyclic low-density parity
check codes, TELSIKS 2003, vol. 2, pp. 417--420.
\href{https://doi.org/10.1109/TELSKS.2003.1246258}{doi:10.1109/TELSKS.2003.1246258}.
\bibitem{Noda} M. Noda, Designing a Self-orthogonal Quasi-cyclic Code with
Extended Minimum Hamming Distance, TURBO-CODING 2006, Munich, 3--7 April 2006.
\href{https://www.vde-verlag.de/proceedings-en/442947110.html}{Publisher record}.
\bibitem{KW} C. A. Kelley and J. L. Walker, LDPC codes from voltage graphs,
ISIT 2008. \href{https://doi.org/10.1109/ISIT.2008.4595095}{doi:10.1109/ISIT.2008.4595095}.
\href{https://digitalcommons.unl.edu/mathfacpub/171/}{Institutional record}.
\bibitem{ABA} R. Asvadi, A. H. Banihashemi and M. Ahmadian-Attari, Lowering the
Error Floor of LDPC Codes Using Cyclic Liftings,
\href{https://arxiv.org/abs/1002.4311}{arXiv:1002.4311v1}, 23 February 2010.
\bibitem{WDW11} J. Wang, L. Dolecek and R. Wesel, LDPC Absorbing Sets, the Null
Space of the Cycle Consistency Matrix, and Tanner's Constructions, ITA 2011.
\href{https://www.seas.ucla.edu/csl/files/publications/JiadongWang_ITA2011.pdf}{Author conference PDF}.
\bibitem{WDW12} J. Wang, L. Dolecek and R. Wesel, The Cycle Consistency Matrix
Approach to Absorbing Sets in Separable Circulant-Based LDPC Codes,
\href{https://arxiv.org/abs/1208.6094v2}{arXiv:1208.6094v2}, 5 December 2012;
first preprint 30 August 2012. IEEE Trans. Inform. Theory 59(4) (2013),
2293--2314.
\bibitem{PP} M. Pra\.{z}mowska and K. Pra\.{z}mowski, Binomial partial Steiner
triple systems containing complete graphs, Graphs and Combinatorics 32 (2016),
2079--2092. \href{https://doi.org/10.1007/s00373-016-1681-3}{doi:10.1007/s00373-016-1681-3}.
\bibitem{AFKL} M. Abreu, M. Funk, V. Kr\v{c}adinac and D. Labbate, Strongly
regular configurations, \href{https://arxiv.org/abs/2104.04880v2}{arXiv:2104.04880v2},
12 September 2021; first preprint 10 April 2021. Designs, Codes and Cryptography
90 (2022), 1881--1897.
\bibitem{AC} A. Abiad and N. Castriota, Spectral and Additive Combinatorial
Methods for Cycles and Absorbing Sets in Lifted-Product Quantum LDPC Codes,
\href{https://arxiv.org/abs/2607.13666v1}{arXiv:2607.13666v1}, 15 July 2026.
Remark 4.18 and Section 6 pose the counting direction; Definition 3.1 and
Section 3.2 give the construction conventions.
\end{thebibliography}
\end{document}
